% GENERATED FILE. Edit canonical lessons, then run npm run generate:books.
% canonical-source-hash: fnv1a64:e3f85bb732e4672e
% canonical-lessons: MW-C154-bolu-hu, MW-C154-bolai-hai, MW-C154-bolo-ho, MW-R154-every-day

\chapter{What I Do Every Day}
\label{ch:mw-a2-what-i-do-every-day}
\hlchaptermodality{154}

\begin{chapteropening}
\textbf{By the end of this chapter:} \emph{I can say what I do, what he does and what you do, every day, with any verb whose stem ends in a consonant.}

Complete MW-R154-every-day: say what you, he and you (politely) do every day.
\end{chapteropening}

\section[\texorpdfstring{\mw{बोलूं} \mw{हूं} (bolū̃ hū̃)}{बोलूं हूं (bolū̃ hū̃)}]{\mw{हूं} \mw{बोलूं} \mw{हूं} — what I do}

\label{lesson:MW-C154-bolu-hu}



\begin{quote}
Say \emph{to speak} and \emph{to hear}. (\textbf{\mw{बोलणो}}, \textbf{\mw{सुणणो}}.) Then \emph{every day}, \textbf{\mw{रोज}}, and \emph{I am here}.
\end{quote}

\begin{grammarlens}[title={Grammar Lens: the stem, -\mw{ूं} and \mw{हूं}}]
Take \textbf{-\mw{णो}} off \textbf{\mw{बोलणो}} and \textbf{\mw{बोल}} is left: the stem. Give it \textbf{-\mw{ूं}}, then end on \textbf{\mw{हूं}}:

\begin{quote}
\textbf{\mw{हूं} \mw{थोड़ु} \mw{बोलूं} \mw{हूं।}}
\end{quote}

\begin{quote}
\emph{hū̃ thoṛu bolū̃ hū̃.} — I speak a little.
\end{quote}

The \textbf{-\mw{ूं}} sounds like \textbf{\mw{हूं}}, \emph{I}: an easy way to hear that the sentence is about you. The same two pieces fit any stem that ends in a consonant: \textbf{\mw{सुणूं} \mw{हूं}}, \textbf{\mw{देखूं} \mw{हूं}}. This is central Marwari as this book writes it; other districts differ.
\end{grammarlens}

\subsection*{Your turn}
Say these aloud:

\begin{itemize}
\raggedright
  \item \textbf{\mw{हूं} \mw{गीत} \mw{सुणूं} \mw{हूं।}} — \emph{hū̃ gīt suṇū̃ hū̃.} — I listen to songs.
  \item \textbf{\mw{हूं} \mw{रोज} \mw{गीत} \mw{सुणूं} \mw{हूं।}} — \emph{hū̃ roj gīt suṇū̃ hū̃.} — I listen to songs every day.
  \item \textbf{\mw{हूं} \mw{थोड़ु} \mw{बोलूं} \mw{हूं।}} — \emph{hū̃ thoṛu bolū̃ hū̃.} — I speak a little.
\end{itemize}

\begin{tcolorbox}[breakable,colback=teal!4,colframe=teal!35!black,title={Before you move on}]
What comes off \textbf{\mw{सुणणो}}? (\textbf{-\mw{णो}}.) What goes on, and what ends the sentence? (\textbf{-\mw{ूं}}, then \textbf{\mw{हूं}}.)
\end{tcolorbox}
\section[\texorpdfstring{\mw{बोलै} \mw{है} (bolai hai)}{बोलै है (bolai hai)}]{\mw{वो} \mw{बोलै} \mw{है} — what he does}

\label{lesson:MW-C154-bolai-hai}



\begin{quote}
Say \emph{I listen to songs}. Then \emph{to read} and \emph{to write}: \textbf{\mw{बांचणो}}, \textbf{\mw{लिखणो}}.
\end{quote}

\begin{grammarlens}[title={Grammar Lens: -\mw{ै} and \mw{है}}]
For \textbf{\mw{वो}}, for \textbf{\mw{राम}}, for anyone you talk about, the stem takes \textbf{-\mw{ै}} and the sentence ends on \textbf{\mw{है}}:

\begin{quote}
\textbf{\mw{वो} \mw{थोड़ु} \mw{बोलै} \mw{है।}}
\end{quote}

\begin{quote}
\emph{vo thoṛu bolai hai.} — He speaks a little.
\end{quote}

\begin{quote}
\textbf{\mw{राम} \mw{रोज} \mw{अखबार} \mw{बांचै} \mw{है।}}
\end{quote}

\begin{quote}
\emph{rām roj akhbār bā̃chai hai.} — Ram reads the newspaper every day.
\end{quote}

Nothing changes for a woman: \textbf{\mw{मां} \mw{गीत} \mw{सुणै} \mw{है।}} (\emph{mā̃ gīt suṇai hai.}), mother listens to songs. Old Marwari wrote the same ending \emph{-aï}: \emph{karaï}, \textquotedblleft{}does\textquotedblright{} (CDIAL 2814).
\end{grammarlens}

\subsection*{Your turn}
Say these aloud:

\begin{itemize}
\raggedright
  \item \textbf{\mw{हूं} \mw{रोज} \mw{अखबार} \mw{बांचूं} \mw{हूं।}} — \emph{hū̃ roj akhbār bā̃chū̃ hū̃.} — I read the newspaper every day.
  \item \textbf{\mw{राम} \mw{रोज} \mw{अखबार} \mw{बांचै} \mw{है।}} — \emph{rām roj akhbār bā̃chai hai.} — Ram reads it every day.
  \item \textbf{\mw{मां} \mw{चिट्ठी} \mw{लिखै} \mw{है।}} — \emph{mā̃ chiṭṭhī likhai hai.} — Mother writes a letter.
\end{itemize}

\begin{tcolorbox}[breakable,colback=teal!4,colframe=teal!35!black,title={Before you move on}]
\emph{I read}? (\textbf{\mw{हूं} \mw{बांचूं} \mw{हूं}}.) \emph{He reads}? (\textbf{\mw{वो} \mw{बांचै} \mw{है}}.)
\end{tcolorbox}
\section[\texorpdfstring{\mw{बोलो} \mw{हो} (bolo ho)}{बोलो हो (bolo ho)}]{\mw{आप} \mw{बोलो} \mw{हो} — what you do}

\label{lesson:MW-C154-bolo-ho}



\begin{quote}
Say \emph{mother listens to songs}. Then \emph{you are here}, and the request \emph{do it}, \textbf{\mw{करो}}.
\end{quote}

\begin{grammarlens}[title={Grammar Lens: -\mw{ो} and \mw{हो}}]
After \textbf{\mw{आप}}, the stem takes \textbf{-\mw{ो}}, the ending of the requests \textbf{\mw{करो}} and \textbf{\mw{लावो}}, and the sentence ends on \textbf{\mw{हो}}:

\begin{quote}
\textbf{\mw{आप} \mw{थोड़ु} \mw{बोलो} \mw{हो।}}
\end{quote}

\begin{quote}
\emph{āp thoṛu bolo ho.} — You speak a little.
\end{quote}

\begin{quote}
\textbf{\mw{आप} \mw{कांई} \mw{करो} \mw{हो}?}
\end{quote}

\begin{quote}
\emph{āp kāĩ karo ho?} — What do you do?
\end{quote}

\textbf{\mw{करो}} alone asks someone to do something. \textbf{\mw{करो} \mw{हो}} after \textbf{\mw{आप}} says what they do.
\end{grammarlens}

\subsection*{Your turn}
Say these aloud:

\begin{itemize}
\raggedright
  \item \textbf{\mw{आप} \mw{गीत} \mw{सुणो} \mw{हो।}} — \emph{āp gīt suṇo ho.} — You listen to songs.
  \item \textbf{\mw{आप} \mw{कांई} \mw{करो} \mw{हो}?} — \emph{āp kāĩ karo ho?} — What do you do?
  \item \textbf{\mw{हूं} \mw{चिट्ठी} \mw{लिखूं} \mw{हूं।}} — \emph{hū̃ chiṭṭhī likhū̃ hū̃.} — I write letters.
\end{itemize}

\begin{tcolorbox}[breakable,colback=teal!4,colframe=teal!35!black,title={Before you move on}]
\emph{I do}? (\textbf{\mw{करूं} \mw{हूं}}: the stem \textbf{\mw{कर}}, \textbf{-\mw{ूं}}, \textbf{\mw{हूं}}.) \emph{He does}? (\textbf{\mw{करै} \mw{है}}.) \emph{You do}, to \textbf{\mw{आप}}? (\textbf{\mw{करो} \mw{हो}}.)
\end{tcolorbox}
\section[(dialogue)]{Every day — a review}

\label{lesson:MW-R154-every-day}



\begin{quote}
Say \emph{I speak a little}.
\end{quote}

\subsection*{Your turn}
Every cell of the present so far, now that each one is taught:

\noindent
\begin{tabularx}{\linewidth}{@{}>{\raggedright\arraybackslash}X>{\raggedright\arraybackslash}X>{\raggedright\arraybackslash}X@{}}
\toprule
\textbf{who} & \textbf{to speak} & \textbf{to read} \\
\midrule
\textbf{\mw{हूं}} & \textbf{\mw{बोलूं} \mw{हूं}} & \textbf{\mw{बांचूं} \mw{हूं}} \\
\textbf{\mw{वो}} & \textbf{\mw{बोलै} \mw{है}} & \textbf{\mw{बांचै} \mw{है}} \\
\textbf{\mw{आप}} & \textbf{\mw{बोलो} \mw{हो}} & \textbf{\mw{बांचो} \mw{हो}} \\
\bottomrule
\end{tabularx}

Say these aloud:

\begin{itemize}
\raggedright
  \item \emph{I listen to songs every day}; \emph{mother listens to songs}
  \item \emph{Ram reads the newspaper every day}; \emph{you read the newspaper}
  \item \emph{what do you do?}, to someone you call \textbf{\mw{आप}}
  \item \emph{I write letters}; \emph{he writes letters}
\end{itemize}

\begin{tcolorbox}[breakable,colback=teal!4,colframe=teal!35!black,title={Before you move on}]
\emph{I speak}? (\textbf{\mw{बोलूं} \mw{हूं}}.) \emph{He speaks}? (\textbf{\mw{बोलै} \mw{है}}.) \emph{You speak}? (\textbf{\mw{बोलो} \mw{हो}}.)
\end{tcolorbox}
